\color{unidarkgreen}{\subsection{Автоматическое ускорение при включении выключателя}}
\color{black}

\begin{enumerate}

\item
В устройстве реализована логика автоматического ускорения (АУ) ДЗ, ТНЗНП, МТЗ и АМТЗ при
включении выключателя.

\item
Функциональная схема АУ представлена на рисунке \ref{pict_au:AA}.

\vspace{2mm}
\begin{figure}[H]
\centering
\includegraphics[width=0.8\textwidth,height=0.8\textheight,keepaspectratio]{\fbpath/08.01 СВ/661.1321 АУ/img/AA.pdf}
\captionof{figure}{Функциональная схема АУ}\label{pict_au:AA}
\end{figure}

\item
Логика АУ вводится по факту включения выключателя и отсутствия симметричного напряжения на одной
из двух систем шин. Действие автоматического ускорения осуществляется от пусков ступеней ДЗ, ТНЗНП, МТЗ
и АМТЗ с регулируемыми независимыми выдержками времени \textbf{«Тср АУ ДЗ»}, \textbf{«Тср АУ ТНЗНП»}, \textbf{«Тср АУ МТЗ»}
и \textbf{«Тср АУ АМТЗ»}. Ввод в работу АУ осуществляется с помощью программного ключа \textbf{«АУ»}. 

\item
Ввод АУ от пуска 2 ступени АМТЗ осуществляется с помощью программного ключа \textbf{«АУ АМТЗ»}.

\item
Пороги срабатывания напряжения прямой последовательности задаются уставками \textbf{«U1 1СШ»} и \textbf{«U1
2СШ»}, напряжения обратной последовательности задаются уставками \textbf{«U2 1СШ»} и \textbf{«U2 2СШ»} и напряжения
нулевой последовательности задаются уставками \textbf{«3U0 1СШ»} и \textbf{«3U0 2СШ»}.

\item
Ввод автоматического ускорения осуществляется по факту включения выключателя (контроль
положения) и отсутствия симметричного напряжения на одной из двух систем шин, на время задаваемое с
помощью уставки \textbf{«Тввод АУ»}.  Контроль отсутствия симметричного напряжения при АУ задается с помощью
программного ключа \textbf{«Контр напряж АУ»}.

\item
Оперативный ввод/вывод АУ производится с помощью БУ (блоков управления функциями) \textit{«Вывод АУ»}.

\item
Уставочные параметры функции АУ представлены в таблице \ref{app_set:au}.

\end{enumerate}